\documentclass[10pt,a4paper]{amsart}
\usepackage[margin=19mm]{geometry}
\usepackage{amsmath,amssymb,mathtools}
\usepackage[scr=rsfs,cal=euler]{mathalfa}
\usepackage{xfrac}
\usepackage[unicode,hidelinks,bookmarks=false]{hyperref}
\newcommand{\ie}{i.e.\ }
\newcommand{\cf}{cf.\ }
\newcommand{\resp}{respectively\ }
\newcommand{\Subsection}{\subsection}
\providecommand{\nobreakdash}{\nobreak\hbox{-}}
\setlength{\emergencystretch}{2em}
\multlinegap0pt
\usepackage{xfrac}
\usepackage{tikz}
\usepackage{dsfont}
\usepackage{bm}
\usepackage{amssymb}
\usepackage{algorithm}\usepackage{algpseudocode}
\usepackage[normalem]{ulem}
\newtheorem{assumption}{Assumption}
\newtheorem{proposition}{Proposition}
\usepackage{cleveref}
\crefname{assumption}{Assumption}{Assumptions}
\crefname{proposition}{Proposition}{Propositions}
\newcommand{\IN}{\mathbb{N}}
\newcommand{\IR}{\mathbb{R}}
\newcommand{\abs}[1]{{\left\vert #1 \right\vert}}
\newcommand{\cA}{\mathcal{A}}
\newcommand{\cG}{\mathcal{G}}
\newcommand{\cP}{\mathcal{P}}
\newcommand{\cW}{\mathcal{W}}
\title{A Multilevel Interacting Particle System Method for the estimation of Failure Probabilities}
\author{Rub\'en Aylwin, Jos\'e Pinto}
\date{Source: arXiv:2608.27275v1 (CC BY 4.0)}
\begin{document}
\maketitle

\noindent\emph{Source and licence.} A Multilevel Interacting Particle System Method for the estimation of Failure Probabilities. Version arXiv:2608.27275v1 (CC BY 4.0), distributed by arXiv under the Creative Commons Attribution 4.0 International licence, http://creativecommons.org/licenses/by/4.0/, as recorded in the arXiv licence metadata for this exact version and shown on the abstract page https://arxiv.org/abs/2608.27275v1. Full licence notice: \texttt{licenses/arxiv-2608.27275-CC-BY-4.0.txt}.\par\smallskip

\noindent\textit{Editorial scope.} Counted unit: the article's proposition prop:setdefi (source0.tex lines 689--701). The article's complete original proof of it is delivered with it (source0.tex lines 702--719, one \texttt{proof} environment of 18 lines). No mathematics is written, repaired or re-derived: the statement and the proof are the article's own lines, in the article's own notation, and no retained line is substituted or re-flowed.  Retained uncounted as the source-local context the proof needs: lines 476--486.  Nothing was omitted from any retained span, so the package carries no omission record.  No retained line contains a citation, so the wrapper carries no bibliography and no bibliography entry.  No retained line contains a figure environment, an image-inclusion command or drawing code, so no image is delivered and the package declares no asset dependency.  Editorial formatting only, in addition: an AMS \texttt{amsart} preamble that restates the article's own declarations and macros used by the retained lines, namely the environments \texttt{assumption}, \texttt{proposition} on separate counters, together with the standard packages those lines need.  The retained environments print in this document as Assumption 1, Proposition 1 in order of appearance, whereas the article prints the same items with the numbering of its own full document.  The complete native source of the article as distributed in the arXiv e-print for this exact version is bundled unchanged as \texttt{sources/208705/source0.tex}.
\begin{assumption}
  \label{as:numap}
  For every $\ell\in\IN_0$ and $y\in\cP$, the computational cost of computing the
  approximation $\cG_\ell(y)=g(u_\ell(y))$ is denoted by $\cW_\ell>0$. The sequence
  $\{\cW_\ell\}_{\ell\in\IN_0}$ is monotonically increasing and there exists a constant
  $\alpha\in(0,1)$ and two exponents $q,\,r>0$ such that
  \begin{align*}
    \abs{\cG(y)-\cG_\ell(y)}\leq C_{\cG}\alpha^{q\ell}\quad\text{and}\quad \cW_\ell \lesssim\alpha^{-r\ell},
  \end{align*}
  for every $\ell\in\IN_0$ and $y\in\cP$, where both $C_{\cG}>0$ and the hidden constant are independent of both $\ell$ and $y$.
\end{assumption}

\begin{proposition}
  \label{prop:setdefi}
  Let \Cref{as:numap} hold, $\{\beta_\ell\}_{\ell=0}^{L-1}\subset\IR^+$ be a finite and non-increasing sequence, and
  define the sets $\{\cA_\ell\}_{\ell=0}^{L-1}$ as follows:
  \begin{align*}
    \cA_\ell:=\left\lbrace y\in\cP\,\Bigg\vert\,\abs{\cG_\ell(y)}\leq C_\cG(\alpha^q+1)\sum\limits_{j=\ell}^{L-1}\alpha^{q j}+\beta_\ell\right\rbrace.
  \end{align*}
  Then, it holds that
  \begin{align*}
    \left\{y\in\cP\,\vert\,\abs{\cG_\ell(y)}\leq C_\cG\alpha^{q\ell}(\alpha^q+1)\right\}\subseteq\cA_\ell\quad\text{and}\quad\cA^{\ell+1}\subseteq\cA^{\ell},
  \end{align*}
  for all $\ell\in\{0:L-1\}$.
\end{proposition}

\begin{proof}
  The first property follows from the fact that $\alpha^{q\ell}\leq\sum_{l=\ell}^{L-1}\alpha^{qj}$ for all $\ell\in\{0:L-1\}$.
  Now pick any $\ell\in\{0:L-2\}$ and consider $y\in\cA_\ell^c$ such that $$\abs{\cG_{\ell}(y)}>C_\cG(\alpha^q+1)\sum_{j=\ell}^{L-1}\alpha^{qj}+\beta_\ell.$$
  Moreover, it follows directly from Assumption \ref{as:numap} that
  \begin{align*}
    \abs{\cG_{\ell}(y)}
    &\leq\abs{\cG_{\ell+1}(y)-\cG_{\ell}(y)}+\abs{\cG_{\ell+1}(y)}\\
    &\leq\abs{\cG_{\ell+1}(y)-\cG(y)}+\abs{\cG(y)-\cG_{\ell}(y)}+\abs{\cG_{\ell+1}(y)}\\
    &\leq C_\cG\alpha^{q\ell}(\alpha^{q}+1)+\abs{\cG_{\ell+1}(y)},
  \end{align*}
  which, for $y \in \mathcal{A}_\ell^c$, implies that
  \begin{align*}
    \abs{\cG_{\ell+1}(y)} &> \beta_{\ell} + C_\cG(\alpha^q+1)\sum\limits_{j=\ell+1}^{L-1}\alpha^{qj}\\
                          &\geq \beta_{\ell+1} + C_\cG(\alpha^q+1)\sum\limits_{j=\ell+1}^{L-1}\alpha^{qj}\implies y\in\cA_{\ell+1}^c,
  \end{align*}
  and we have shown that $\cA_{\ell}^c\subseteq\cA_{\ell+1}^c$, which yields the required property after taking
  the complement.
\end{proof}

\end{document}
